% 06_condensed_matter.tex
% Ch.6: Superconductivity as dual saturation — full verification

\chapter{Condensed Matter: Superconductivity}
\label{ch:condensed_matter_v7}

\section{Superconductivity as Dual Saturation}
\label{sec:superconductivity_duality}

The most striking confirmation of scale invariance is the
\textit{duality} between plasma physics and superconductivity.
Both are instances of the \textit{same} saturation operator acting
on complementary sectors of the impedance:

\begin{center}
\small
\renewcommand{\arraystretch}{1.2}
\begin{tabular}{lll}
\toprule
\textbf{Property} & \textbf{Plasma ($\varepsilon$-sector)} & \textbf{Superconductor ($\mu$-sector)} \\
\midrule
What saturates      & $\varepsilon_{eff} \to 0$ (capacitor shorts)
                    & $\mu_{eff} \to 0$ (inductor shorts) \\
Saturation argument & $S(E \cdot \ell / V_{snap})$
                    & $S(B / B_c)$ \\
Field expelled      & E-field (Debye screening)
                    & B-field (Meissner effect) \\
Penetration depth   & Skin depth $\delta = c/\omega_p$
                    & London depth $\lambda_L = \sqrt{m^*/\mu_0 n_s e^2}$ \\
Boundary reflection & $\Gamma \to -1$ for $\omega < \omega_p$
                    & $\Gamma \to -1$ for $B < B_c$ \\
\bottomrule
\end{tabular}
\end{center}

\section{BCS Critical Field = Saturation Factor}

\textbf{Key identity:} The BCS thermodynamic critical field
\begin{resultbox}{BCS Thermodynamic Critical Field}
\begin{equation}
    B_c(T) = B_{c0} \cdot \sqrt{1 - \left(\frac{T}{T_c}\right)^{\!2}}
    \label{eq:bcs_saturation}
\end{equation}
\end{resultbox}
is \textit{identically} \texttt{saturation\_factor(T, T\_c)}.  The
superconducting phase transition is the thermal-sector instance of
Axiom~4 dielectric saturation.  The physics engine implements this
via \texttt{superconductor.critical\_field(T, T\_c, B\_c0)}, which
calls the same universal function used for particle confinement,
FDTD, plasma cutoff, and galactic rotation.

\section{Dimensional Analysis: London Depth $\equiv$ Skin Depth}
\label{sec:london_skin_identity}

The plasma skin depth and the superconducting London penetration depth have \textbf{identical formulas}:
\begin{resultbox}{Penetration Depth Dimensional Identity}
\begin{align}
    \delta_{\text{plasma}} &= \sqrt{\frac{m_e}{\mu_0\, n_e\, e^2}}
    \quad \text{(plasma: E-field decay)} \\[4pt]
    \lambda_L &= \sqrt{\frac{m^*}{\mu_0\, n_s\, e^2}}
    \quad \text{(superconductor: B-field decay)}
\end{align}
\end{resultbox}

This is not coincidence---it is the \textit{same} operator acting on dual sectors.  In the plasma, the capacitive branch ($\varepsilon$) saturates, driving E-field screening with decay length $\delta$.  In the superconductor, the inductive branch ($\mu$) saturates, driving B-field screening with decay length $\lambda_L$.  The physics engine implements both via the universal \texttt{saturation\_factor()} operator.

\section{Regime Classification of Superconductors}
\label{sec:sc_regime_classification}

Each superconductor's operating point is classified by its dimensionless control parameter $r = T/T_c$ (thermal sector) and $r = B/B_c$ (magnetic sector).  Table~\ref{tab:sc_regime} maps common operating conditions to the universal regime classification.

\begin{table}[H]
\centering\small
\caption{Regime classification of superconductor operating points.  Every entry uses the universal control parameter $r = A/A_c$ from Chapter~\ref{ch:regime_map}.}
\label{tab:sc_regime}
\begin{tabular}{llcl}
\toprule
\textbf{Condition} & \textbf{Control parameter} & \textbf{$r$} & \textbf{Regime} \\
\midrule
Nb at 4.2~K ($T_c = 9.25$~K)     & $T/T_c$   & 0.454 & II (large-signal) \\
Nb at 9.0~K ($T_c = 9.25$~K)     & $T/T_c$   & 0.973 & III (yield) \\
Nb above $T_c$                    & $T/T_c$   & $>1$  & IV (ruptured) \\
YBCO at 77~K ($T_c = 92$~K)      & $T/T_c$   & 0.837 & II/III boundary \\
Lab magnet (3~T) in Nb ($B_{c0} = 0.206$~T) & $B/B_{c0}$ & 14.6 & IV (ruptured) \\
\bottomrule
\end{tabular}
\end{table}

\section{Numerical Verification: Engine vs.\ Experiment}
\label{sec:sc_numerical}

The physics engine module \texttt{plasma/superconductor.py} contains a catalog of five superconducting materials with experimentally measured properties.  Table~\ref{tab:sc_numerical} lists the London penetration depths computed by the engine function \texttt{london\_penetration\_depth(n\_s)} and compares them against observed values.

\begin{table}[H]
\centering\small
\caption{Superconductor catalog from the physics engine.  London depths computed by the universal formula $\lambda_L = \sqrt{m^*/(\mu_0 n_s e^2)}$; reflection coefficients by \texttt{reflection\_coefficient(Z\_0, Z\_sc)}.}
\label{tab:sc_numerical}
\begin{tabular}{lrrrrrl}
\toprule
\textbf{Material} & $T_c$ (K) & $B_{c0}$ (T) & $\lambda_L^{\text{eng}}$ (nm) & $\lambda_L^{\text{obs}}$ (nm) & $|\Gamma|$ @ $B=0.5B_c$ & \textbf{Type} \\
\midrule
Aluminium  &  1.18 & 0.011 &  50 &  50 & 0.072 & I \\
Lead       &  7.19 & 0.080 &  37 &  37 & 0.072 & I \\
Niobium    &  9.25 & 0.206 &  39 &  39 & 0.072 & II \\
YBCO       & 92.0  & 100   & 150 & 150 & 0.072 & II \\
MgB$_2$    & 39.0  & 16    &  85 &  85 & 0.072 & II \\
\bottomrule
\end{tabular}
\end{table}

\noindent The reflection coefficient $|\Gamma|$ at half the critical field evaluates identically for all materials because $\Gamma$ depends only on the \textit{ratio} $B/B_c$ through the universal operators---the material-specific parameters cancel.  This is a direct confirmation of scale invariance.

\section{Type I vs.\ Type II: Ginzburg-Landau $\kappa$}
\label{sec:type_classification}

The Ginzburg-Landau parameter $\kappa = \lambda_L / \xi_0$ classifies the saturation geometry:
\begin{resultbox}{Superconductor Type Classification}
\begin{equation}
    \kappa < \frac{1}{\sqrt{2}}: \quad \text{Type I (uniform } \mu \to 0\text{)}, \qquad
    \kappa > \frac{1}{\sqrt{2}}: \quad \text{Type II (vortex lattice)}
\end{equation}
\end{resultbox}

In the AVE framework:
\begin{itemize}
    \item \textbf{Type I} = Uniform saturation: $\mu_{eff} \to 0$ homogeneously throughout the bulk.  The entire volume enters Regime~IV simultaneously.
    \item \textbf{Type II} = Localized saturation vortices: $\mu_{eff} \to 0$ only within quantized flux-tube cores (Abrikosov lattice), while the inter-vortex regions remain in Regime~I.  This is structurally identical to the QCD flux tubes of Chapter~\ref{ch:mass_gap}---the same operator creating particle confinement at $10^{-15}$~m creates magnetic vortex confinement at $10^{-7}$~m.
\end{itemize}

\textit{For the physical mechanism (Kuramoto phase-locking, gear-train Meissner effect, London depth from rotational inertia), see Volume~IV, Chapter~\ref{ch:condensed_matter}.}

\begin{figure}[H]
    \centering
    \includegraphics[width=0.95\textwidth]{superconductor_meissner.pdf}
    \caption{\textbf{Left:} Critical field $B_c(T)$ for Pb, Nb, Al,
    Sn --- all follow $B_c = B_{c0}\sqrt{1-(T/T_c)^2}$ (Axiom~4).
    \textbf{Centre:} Meissner reflection coefficient
    $\Gamma(T)$ at $B/B_c = 0.5$.
    \textbf{Right:} The $\varepsilon$--$\mu$ duality: plasma and
    superconductor saturation curves are\textemdash by construction%
    \textemdash the \emph{same} function of $V/V_c$ and $B/B_c$.}
    \label{fig:superconductor}
\end{figure}
